% Part: counterfactuals
% Chapter: introduction
% Section: counterfactuals

\documentclass[../../../include/open-logic-section]{subfiles}

\begin{document}

\olfileid{cnt}{int}{cnt}

\olsection{خلافِ واقعہ شرطیے}

انگریزی وچ ’جے \dots تاں \dots‘ دی اک بڑی عام تے ضروری شکل
\emph{to be} دے ماضی والے فرضیہ صیغے (past subjunctive) نال بندی اے۔
ایہدی شکل اے ’جے ایہ گل ہوندی کہ \dots تاں ایہ گل ہوندی کہ \dots‘۔
ایہو جیہے شرطیے دا مقدم عام طور اُتے جھوٹا ہوندا اے، یعنی حقیقت دے اُلٹ۔
ایس لئی ایہناں نوں \emph{خلافِ واقعہ شرطیے} آکھدے نیں۔ \emph{to be}
دا فرضیہ صیغہ ورتن کارن ایہناں نوں \emph{فرضیہ شرطیے} وی آکھدے نیں۔
ایہ \emph{خبری} شرطیاں توں وکھرے نیں، جیہناں دی شکل اے ’جے ایہ گل اے
کہ \dots تاں ایہ گل اے کہ \dots‘۔ خلافِ واقعہ تے خبری شرطیاں دے سچ
ہون دیاں شرطاں وکھریاں نیں۔ ایڈمز دی مشہور مثال ویکھو:
\begin{quote}
  جے اوسوالڈ نے کینیڈی نوں نہیں ماریا، تاں کسے ہور نے ماریا اے۔

  جے اوسوالڈ نے کینیڈی نوں نہ ماریا ہوندا، تاں کسے ہور نے مار دینا سی۔
\end{quote}
پہلا خبری اے، دوجا خلافِ واقعہ۔ پہلا صاف طور اُتے سچ اے۔ سانوں پتا اے
کہ صدر جان ایف کینیڈی نوں \emph{کسے} نے ماریا سی، تے جے (وارن رپورٹ دے اُلٹ)
اوہ بندہ لی ہاروی اوسوالڈ نہیں سی، تاں کینیڈی نوں کسے ہور نے ماریا سی۔
دوجا کجھ ہور آکھدا اے۔ ایہ دعویٰ کردا اے کہ جے اوسوالڈ نے کینیڈی نوں نہ
ماریا ہوندا، یعنی جے ڈیلس دی گولی باری ٹل گئی ہوندی یا ناکام رہی ہوندی،
تاں اگوں تاریخ ایویں ودھدی کہ قتل دی اک ہور کوشش کامیاب ہو جاندی۔
ایہدے سچ ہون لئی ضروری اے کہ طاقتور دھڑیاں نے کینیڈی دی موت یقینی بناؤن
لئی سازش کیتی ہوندی (جیویں کینیڈی دے قتل دی سازش منّن والے کئی لوک
سمجھدے نیں)۔

ایہ بحث ہجے جاری اے کہ مادی شرطیہ \emph{خبری} شرطیے دا مطلب ٹھیک طرح
دیندا اے یا نہیں۔ خاص کر کے بحث ایہ اے کہ مادی شرطیے دے معمّیاں دی
’سمجھاؤنی‘ ایویں ہو سکدی اے یا نہیں کہ اوہ انگریزی دے خبری شرطیاں دے
سچ ہون دیاں شرطاں وی ٹھیک دیندا ہووے۔ البتہ ایس گل اُتے اختلاف نہیں کہ
مادی شرطیاں نال خلافِ واقعہ شرطیے ٹھیک علامتاں وچ نہیں لکھے جا سکدے۔
ایہ گل ایس توں صاف اے کہ بھانویں خلافِ واقعہ شرطیاں دے مقدم عام طور اُتے
جھوٹے نیں، جھوٹے مقدم والے سارے خلافِ واقعہ شرطیے سچ نہیں۔ مثلاً جے تسی
وارن رپورٹ منّ دے او تے کینیڈی دے قتل دی کوئی سازش نہیں سی، تاں ایڈمز دا
خلافِ واقعہ شرطیہ ایہدی مثال اے۔

سبب بارے دلیل وچ خلافِ واقعہ شرطیاں دا وڈا کم اے۔ ایہناں دی اک
اہم ورتوں سبب تے اثر دا تعلق بیان کرنا اے۔ مثلاً تیلی رگڑن نال اوہ بل پَیندی اے،
تے ایہ گل ایویں آکھ سکدے آں، ’جے ایہ تیلی رگڑی جاندی، تاں بل پَیندی۔‘
مادی، تے عام طور اُتے خبری، شرطیاں نال ایہ گل نہیں آکھی جا سکدی۔
’تیلی رگڑی جاندی اے $\lif$ تیلی بلدی اے‘ سچ اے جے تیلی کدی رگڑی
ای نہ جاوے، بھانویں رگڑن اُتے جو مرضی ہوندا۔ ہور وی ماڑی گل، ’تیلی رگڑی
جاندی اے $\lif$ تیلی پھُلّاں دے گلدستے وچ بدل جاندی اے‘ وی سچ اے جے
تیلی کدی رگڑی ای نہ جاوے۔ پر جے اوہ رگڑی جاندی، تاں اوہ یقیناً پھُلّاں
دا گلدستہ نہ بندی۔

خلافِ واقعہ شرطیاں دی صحیح منطق کیہڑی اے، ایس اُتے ہجے بحث اے۔
سٹالنیکر تے لوئس نے ایہناں شرطیاں دی اک بااثر تشریح پیش کیتی۔
اوہناں مطابق خلافِ واقعہ شرطیہ ’جے ایہ گل ہوندی کہ~$S$، تاں ایہ گل
ہوندی کہ~$T$‘ اُدوں تے صرف اُدوں سچ اے جدوں $T$ اوس خلافِ واقعہ حالت
(’ممکن جہان‘) وچ سچ ہووے جیہڑی اصل جہان دی حالت دے سبھ توں نیڑے اے
تے جتھے~$S$ سچ اے۔ ایہنوں ’وجودی‘ پڑتال آکھدے نیں، کیوں جے ایہ ممکن
جہاناں دی وجودیات دا حوالہ دیندی اے۔ ہور پڑتالاں مشروط احتمالاں یا
یقین بدلّن دے نظریے ورتدیاں نیں۔ خلافِ واقعہ شرطیاں لئی کئی وکھ وکھ
منطقاں پیش کیتیاں گئیاں نیں۔ ایہناں دی کوئی اکّو لوئس--سٹالنیکر منطق وی
نہیں۔ بھانویں سٹالنیکر تے لوئس دوناں دیاں تشریحاں سبھ توں نیڑے ممکن جہاناں
دے خیال اُتے ٹکیاں نیں، اوہناں دے مفروضیاں توں وکھ وکھ معتبر استنباط
نکلدے نیں۔

\end{document}
